% ---------------------------------------------------------------------------
% Section: The held rough error and the closure beyond the lazy regime
%
% Drop-in section for the paper. Conventions of Part III: Pi = I, c_p = 1,
% gradient flow, E = \hat Y - Y. Rewritten 2026-10-12 around the roughness
% function of Part I that the theory singles out, the Krylov functional; the
% flow of the section is unchanged. Cross-references to the main text are
% by label where the main text has one (thm:stepsup, thm:stepkrylov,
% thm:slow, prop:rates, cor:budget, thm:path, thm:time0, cor:mintime,
% thm:duhamel, prop:ell, lem:drift) and by hard-coded number otherwise
% (Theorem 7(a) = energy identity, Proposition 15(b) = rate of sigma_R,
% Proposition 20 = analyticity radius, Corollary 6 = comparison principle);
% replace with \ref{...} to your labels. The theorem environments used are
% proposition, theorem, corollary, definition, remark, proof.
% ---------------------------------------------------------------------------

\section{The held rough error and the closure beyond the lazy regime}
\label{sec:held-rough}

Proposition~20 says where every expansion about $\theta_0$ stops: a rough
direction is fitted only after a displacement of order the analyticity radius
$R$, so no series in $\theta-\theta_0$ describes the stretch of training in
which the rough part of the error is removed. This section changes the
expansion variable. The one-time theory of Part~I bounds the eigenvalue of the
kernel on a rough direction by the peak feature roughness,
$\lambda_k R(v_k)^2 \le \sigma_R^2$ (Corollary~\ref{cor:budget}), and
Proposition~15 gives the exact rate of $\sigma_R$ along training. Together
they say that the part of the error that a roughness function counts as rough
can lose energy only through one scalar function of time, $\sigma_R(t)$; and
the a~priori bound on the growth of $\sigma_R$ then gives a floor on the rough
part of the error that holds at every time, with no ball and no horizon.

One point governs everything below and is the reason the section is written
as it is. A roughness function is a dual variable for a lower bound: every
seminorm gives one, and the supremum over all of them is attained by a single
rank-one functional built from the current error and the current kernel, the
slow part $E_\gamma=\gamma(K+\gamma I)^{-1}E$ normalised,
$\varphi_\gamma=E_\gamma/\|E_\gamma\|$, which is a polynomial in $K$ applied
to $E$ and is computed by the Krylov series
(Theorems~\ref{thm:stepsup}, \ref{thm:stepkrylov} and~\ref{thm:slow}(c)). For
it, every inequality in the floor is an equality in its signed form. A
roughness function fixed by the inputs alone, such as the cluster seminorm or
the manifold high-pass, is a local complexity function: it makes the roughness
of the labels computable before training and gives a weaker bound. In
practice the only roughness function worth using is therefore the Krylov
functional of the current state, and every statement below is made for it
first; the fixed-seminorm statements are its corollaries. An earlier draft
fixed a Laplacian high-pass throughout, measured its Rayleigh-quotient
closure to fail along training, and replaced the closure by a model; that
failure is an artefact of the fixed projector, as \S\ref{ssec:hr-closure}
explains, and the model is gone.

Throughout, $R(u)=\|Pu\|$ is a projector seminorm, with $P$ a symmetric
projector on sampled functions, so that $C_R=1$ and
$\sigma_R(\theta)^2=\lambda_{\max}\big(PK(\theta)P\big)$. The rank-one
functional is of this form with $P=\varphi\varphi^{\!\top}$ for a unit sampled
function $\varphi$, for which $\sigma_R=\|\varphi\|_K$. Write
\begin{equation}
E_r = PE, \qquad E_s = (I-P)E, \qquad \sigma(t)=\sigma_R(\theta(t)),
\label{eq:split-rs}
\end{equation}
for the rough and the smooth part of the error under $R$ and for the peak
feature roughness along the path. Since $Y$ is fixed, $\dot E_r = P\dot E$,
and $\|E_r(0)\| = R(E_0) \ge R(Y)-R(\hat Y_0)$.

\subsection{The rough part loses energy only through $\sigma_R$}

\begin{proposition}[Energy of the rough part]
\label{prop:rough-energy}
Let $N_\theta$ be continuously differentiable in $\theta$ and let $\theta(t)$
be the gradient-flow path. Then, with the notation of
Eq.~\eqref{eq:split-rs}:
\begin{enumerate}
\item[(a)] Exactly,
\begin{equation}
\frac{d}{dt}\,\|E_r\|^2 = -2\,\langle E_r, K E_r\rangle - 2\,\langle E_r, K E_s\rangle,
\qquad
0 \le \langle E_r, K E_r\rangle \le \sigma^2\,\|E_r\|^2,
\qquad
|\langle E_r, K E_s\rangle| \le \sigma\,\lambda_1^{1/2}\,\|E_r\|\,\|E_s\|.
\label{eq:rough-rate}
\end{equation}
\item[(b)] For a unit sampled function $\varphi$ and $P=\varphi\varphi^{\!\top}$,
the first inequality in (a) is an equality,
$\langle E_r,KE_r\rangle=\sigma^2\|E_r\|^2$ with $\sigma=\|\varphi\|_K$, and
the signed component obeys, exactly and for every update rule,
\begin{equation}
\frac{d}{dt}\langle\varphi,E\rangle = -\langle\varphi,KE\rangle,\qquad
\Big|\frac{d}{dt}\langle\varphi,E\rangle\Big| \le \|\varphi\|_{K}\,|\dot\theta| = \sigma\,|\dot\theta| .
\label{eq:signed-rate}
\end{equation}
\item[(c)] The rate at which the rough part can lose energy is bounded by the
peak feature roughness alone once the smooth part is gone:
$-\frac{d}{dt}\|E_r\|^2 \le 2\sigma^2\|E_r\|^2 + 2\sigma\lambda_1^{1/2}\|E_r\|\,\|E_s\|$.
\end{enumerate}
\end{proposition}

\begin{proof}
(a) $\frac{d}{dt}\frac12\|E_r\|^2 = \langle PE, P\dot E\rangle = \langle E_r,\dot E\rangle = -\langle E_r, KE\rangle$
by Theorem~7(a), and $E = E_r+E_s$. $K$ is positive semidefinite, so
$\langle E_r,KE_r\rangle\ge0$. For a unit $\psi$ in the range of $P$,
$R(\psi)=\|\psi\|=1$, so $\|\psi\|_K\le\sigma_R$ by the definition of
$\sigma_R$; hence $\langle E_r,KE_r\rangle=\|E_r\|_K^2\le\sigma^2\|E_r\|^2$.
For the cross term, $\langle a,Kb\rangle = g_a\cdot g_b$, so
$|\langle E_r,KE_s\rangle|\le \|E_r\|_K\|E_s\|_K\le \sigma\|E_r\|\cdot\lambda_1^{1/2}\|E_s\|$.
(b) $E_r=\langle\varphi,E\rangle\varphi$ and $PKP$ has the single nonzero
eigenvalue $\langle\varphi,K\varphi\rangle=\|\varphi\|_K^2$. The rate identity
is $\dot E=\Psi\dot\theta$ paired with $\varphi$,
$\langle\varphi,\Psi\dot\theta\rangle=g_\varphi\cdot\dot\theta$ with
$|g_\varphi|=\|\varphi\|_K$ (Proposition~\ref{prop:rates}(a)); under gradient
flow $\dot\theta=-g_E$ gives the first form. (c) is (a) with the two bounds.
\end{proof}

Part (a) is the energy identity of Theorem~7(a) projected on the rough
subspace. For a general projector it has two terms, the Rayleigh quotient on
the rough subspace, capped by $\sigma^2$, and the coupling to the smooth part,
capped by $\sigma\lambda_1^{1/2}$. Part (b) is what the Krylov functional
changes. The cap on the Rayleigh quotient is attained, so no closure constant
appears; and the component $\langle\varphi,E\rangle$ obeys a rate bound with
no coupling term at all, because a rank-one seminorm measures a signed
component and the sign carries the information that the energy form throws
away. The coupling term of (a) is not small for this functional: on the
reference problem it is $4$ to $12$ times the diagonal term from $t=0$ on
[n=200], so an energy bound for $\|E_r\|^2$ built from (a) is empty for it,
while the signed bound of (b) is exact. Everything below therefore integrates
Eq.~\eqref{eq:signed-rate}, not Eq.~\eqref{eq:rough-rate}.

\subsection{The a priori floor on the rough part}

Integrated along the path, Eq.~\eqref{eq:signed-rate} charges the component
of the error along $\varphi$ by the cost of that direction,
$\mathrm{cost}_t(\varphi)=\int_0^t\|\varphi\|_{K(s)}|\dot\theta(s)|\,ds$
(Theorem~\ref{thm:path}). The a~priori form replaces the kernel along the path
by its value at the start plus the drift, which Lemma~\ref{lem:drift} bounds
by the second derivative of the network over the ball the path stays in:
with $\ell(t)=\int_0^t|\dot\theta|$ the path length and $\bar H_2(\ell)$ the
largest second-derivative size over the ball of radius $\ell$ about
$\theta_0$, $\|\varphi\|_{K(s)}\le\|\varphi\|_{K_0}+\bar H_2(\ell)\,\ell(s)$ for
every unit $\varphi$ (Proposition~15(b) in its path form).

\begin{theorem}[The rough part is held: a priori form]
\label{thm:rough-apriori}
Let $N_\theta$ be twice continuously differentiable in $\theta$ and let
$\theta(t)$ be the gradient-flow path, with $\ell=\ell(t)$. Then
\begin{equation}
\|E(t)\| \;\ge\; F_0(\ell) - \tfrac12\,\bar H_2(\ell)\,\ell^2,
\qquad
F_0(\ell)=\max_{\|\varphi\|\le1}\Big(\langle\varphi,E_0\rangle-\ell\,\|\varphi\|_{K_0}\Big),
\label{eq:rough-floor-apriori}
\end{equation}
where $F_0(\ell)$ is the smallest error of the first-order expansion about
$\theta_0$ of a step of length $\ell$, attained by the Krylov functional of
$(K_0,E_0)$ at the rate $\gamma$ with $\Upsilon_0(\gamma)=\ell^2$, and
computed by the Krylov series (Theorems~\ref{thm:stepsup} and~\ref{thm:stepkrylov};
this is Theorem~\ref{thm:time0}). For every projector seminorm $R$ with
$\sigma_0=\sigma_R(\theta_0)$,
\begin{equation}
\|E(t)\| \;\ge\; R(Y)-R(\hat Y_0)-\ell\,\sigma_0-\tfrac12\,\bar H_2(\ell)\,\ell^2 ,
\label{eq:rough-floor-R}
\end{equation}
and under gradient flow $\ell(t)\le\sqrt{t/2}\,\|E_0\|$ (Proposition~\ref{prop:ell}).
\end{theorem}

\begin{proof}
For a unit $\varphi$, Eq.~\eqref{eq:signed-rate} integrated gives
$\langle\varphi,E(t)\rangle\ge\langle\varphi,E_0\rangle-\mathrm{cost}_t(\varphi)$,
and $\|E(t)\|\ge\langle\varphi,E(t)\rangle$. Insert the drift bound,
$\mathrm{cost}_t(\varphi)\le\ell\|\varphi\|_{K_0}+\bar H_2(\ell)\int_0^t\ell(s)\,d\ell(s)
=\ell\|\varphi\|_{K_0}+\tfrac12\bar H_2(\ell)\ell^2$, and take the maximum over
$\varphi$. For Eq.~\eqref{eq:rough-floor-R}, take $\varphi$ the unit vector in
the range of $P$ that attains $R(E_0)=\langle\varphi,E_0\rangle$, for which
$\|\varphi\|_{K_0}\le\sigma_0$, and use $R(E_0)\ge R(Y)-R(\hat Y_0)$. The path
length bound is Cauchy--Schwarz on $\ell=\int\|E\|_K$ with the energy
identity.
\end{proof}

Three features distinguish the theorem. It holds at every time with no
horizon and no analyticity, since nothing is expanded. Its first term is
exact in the roughness function: $F_0(\ell)$ is the supremum of the
fixed-seminorm floors Eq.~\eqref{eq:rough-floor-R} over all $R$, so no
seminorm has to be chosen and none does better. And the network enters only
through $K_0$, in $F_0$, and through $\bar H_2$: the kernel at the start says
what a path of length $\ell$ can remove to first order, and the second
derivative says how much more it can remove by moving the kernel. The
earlier draft stated this floor for a fixed projector in energy form, with
$\sigma$ bounded by $\sigma_0+as$ and the decay charged exponentially; that
form carries the coupling term of Proposition~\ref{prop:rough-energy}(a) and
is empty for the functional that matters, so it is withdrawn in favour of
Eq.~\eqref{eq:rough-floor-apriori}.

\begin{corollary}[Minimum path length and training time to remove the rough part]
\label{cor:min-time-rough}
Under the hypotheses of Theorem~\ref{thm:rough-apriori}, let
$\ell_\epsilon^\star=\inf\{\ell\ge0: F_0(\ell)-\tfrac12\bar H_2(\ell)\ell^2\le\epsilon\}$.
If $\|E(t)\|\le\epsilon<\|E_0\|$, then $\ell(t)\ge\ell_\epsilon^\star$ and
\begin{equation}
t \;\ge\; \frac{2\,(\ell_\epsilon^\star)^2}{\|E_0\|^2-\epsilon^2} .
\label{eq:min-time-rough}
\end{equation}
For every projector seminorm, $\ell_\epsilon^\star\ge\ell_\epsilon^R$, the
smallest $\ell$ with $\ell\sigma_0+\tfrac12\bar H_2(\ell)\ell^2\ge R(Y)-R(\hat Y_0)-\epsilon$.
\end{corollary}

\begin{proof}
Corollary~\ref{cor:mintime}: the right side of
Eq.~\eqref{eq:rough-floor-apriori} is nonincreasing in $\ell$, and
$\ell(t)^2\le\tfrac t2(\|E_0\|^2-\epsilon^2)$ at the first time the error
reaches $\epsilon$.
\end{proof}

Eq.~\eqref{eq:min-time-rough} is the a~priori ignition time of the rough part
of the labels: the path must be long enough for the first-order term to pay
for the roughness the labels demand, or the second-derivative term must have
moved the kernel, and under gradient flow the path length grows at most as
$\sqrt{t}$. The Krylov form $\ell_\epsilon^\star$ is never worse than any
seminorm's $\ell_\epsilon^R$.

\subsection{The roughness spectrum and the moving cutoff}
\label{ssec:hr-cutoff}

The floor of Theorem~\ref{thm:rough-apriori} is one number. Read across the
rate $\gamma$ of the slow part, it locates where the held error sits. For a
sampled function $u$ write $u_\gamma=\gamma(K_0+\gamma I)^{-1}u$ for its slow
part under the initial kernel; $\gamma\mapsto\|u_\gamma\|$ is nonincreasing,
equal to $\|u\|$ as $\gamma\to\infty$ and to the part of $u$ outside the range
of $K_0$ as $\gamma\to0$.
\begin{equation}
\rho(\gamma) = \|E_{0,\gamma}\|
\label{eq:rough-spectrum}
\end{equation}
is the roughness spectrum of the initial error in the coordinates of the
kernel: the energy of $E_0$ that the initial kernel removes at rate $\gamma$
or slower. For a target with $\hat Y_0$ small it is the roughness spectrum of
the labels, $\|Y_\gamma\|$, the dataset as the kernel sees it
(Part~I, the dataset in the coordinates of the kernel); it is fixed by the
dataset and the initial kernel and computed by the Krylov series.

\begin{definition}[Cutoff rate]
\label{def:cutoff}
For $t\ge0$, the cutoff rate is
\begin{equation}
\gamma^\star(t) = \sup\Big\{\gamma>0 : \ell(t)\,\tfrac{\sqrt\gamma}{2}+\tfrac12\,\bar H_2(\ell(t))\,\ell(t)^2 \le \tfrac12\,\rho(\gamma)\Big\},
\label{eq:cutoff}
\end{equation}
the fastest rate at which, at time $t$, the a~priori bound allows the slow
part of the initial error to have lost at most half its size.
\end{definition}

Since $\rho$ is nonincreasing in $\gamma$ and $\ell$ nondecreasing in $t$,
$\gamma^\star(t)$ is nonincreasing in $t$: the cutoff moves down through the
rates as training proceeds, and it moves even when $\sigma_R(\theta_0)=0$ on
a direction, through the $\bar H_2\ell^2$ term, which is feature learning.

\begin{proposition}[The error at time $t$ is at least the energy above the cutoff]
\label{prop:cutoff-floor}
Under the hypotheses of Theorem~\ref{thm:rough-apriori}, for every $t$ and
every $\gamma>0$,
\begin{equation}
\|E(t)\| \;\ge\; \rho(\gamma)-\ell(t)\,\tfrac{\sqrt\gamma}{2}-\tfrac12\,\bar H_2(\ell(t))\,\ell(t)^2,
\qquad\text{so}\qquad
\|E(t)\|\ge\tfrac12\,\rho(\gamma^\star(t)) .
\label{eq:cutoff-floor}
\end{equation}
\end{proposition}

\begin{proof}
Take $\varphi=\varphi_\gamma=E_{0,\gamma}/\|E_{0,\gamma}\|$ in the proof of
Theorem~\ref{thm:rough-apriori}. In the eigenbasis of $K_0$ with weights
$w_k=\gamma/(\lambda_k+\gamma)$, $\langle\varphi_\gamma,E_0\rangle=\sum_kw_k\hat e_k^2/\|E_{0,\gamma}\|\ge\sum_kw_k^2\hat e_k^2/\|E_{0,\gamma}\|=\rho(\gamma)$,
and $\|\varphi_\gamma\|_{K_0}^2=\sum_k\lambda_kw_k^2\hat e_k^2/\sum_kw_k^2\hat e_k^2\le\max_\lambda\lambda\gamma^2/(\lambda+\gamma)^2=\gamma/4$.
\end{proof}

In words: at time $t$, the error is at least the part of the initial error
that the initial kernel removes more slowly than the cutoff rate, and the
cutoff is set by the path length and the roughness capacity of the network,
$\sqrt\gamma/2$ per unit length along the slow part at rate $\gamma$ and
$\bar H_2$ for the motion of the kernel, and by nothing else. Without feature
learning ($\bar H_2=0$) the cutoff is $\gamma^\star$ with
$\ell\sqrt{\gamma^\star}\approx\rho(\gamma^\star)$, which is the spectral
bias of the lazy regime stated in the kernel's own coordinates; with feature
learning the cutoff falls faster, and Corollary~\ref{cor:min-time-rough} says
how much faster it can fall at most. For labels whose roughness spectrum
$\rho$ has a step at some rate $\gamma_Y$, the error is held near
$\rho(\gamma_Y)$ until $\gamma^\star(t)$ passes below $\gamma_Y$. The chain
section replaces the a~priori charge $\ell\sqrt\gamma/2+\tfrac12\bar H_2\ell^2$
by the actual lifted spectrum along the run: its cutoff bound is this
proposition with the weight $\gamma/(\lambda+\gamma)$ hardened to the
indicator of $\lambda_k(T)T<1$ and the slow set transported along the
eigenvector flow.

\subsection{The closure, revisited}
\label{ssec:hr-closure}
The statements above bound the rough part from below. An earlier draft tried
to close them from above with the hypothesis that the kernel's Rayleigh
quotient on the rough residual stays a fixed fraction $c$ of $\sigma_R^2$,
$\langle E_r,KE_r\rangle\ge c\,\sigma_R^2\|E_r\|^2$, measured it with a fixed
Laplacian high-pass, found $c$ to decay from $0.02$--$0.09$ to $0.001$--$0.005$
along training while $\sigma_R^2$ was lifted $6$--$50\times$ [n=200], and
replaced the closure by a reduced model. That measurement is an artefact of
holding the projector fixed: the lift raises the top of the fixed rough
subspace while the residual moves to the slowest directions of that
subspace, so $c(t)\approx\lambda_{\mathrm{eff}}(t)/\sigma_R(t)^2\to0$. For the
Krylov functional of the current state, $c=1$ identically
(Proposition~\ref{prop:rough-energy}(b)), and the closure hypothesis is not
needed. What is genuinely open is not a closure but the cost: the floor along
the path charges every direction by $\int\|\varphi\|_{K(s)}|\dot\theta|$, and
along the rough directions $\|\varphi\|_{K(s)}$ is lifted by orders of
magnitude, which is feature learning; the a~priori bound on that lift through
$\bar H_2$ is a worst case. The two-sided statement that replaces the band is
the cutoff bound of Section~\ref{sec:chain}: the energy of the initial error
on the transported slow set, with the actual lifted spectrum in place of the
a~priori charge and the cross-front rotation as the named remainder, and the
onset of the lift given by a convergent expansion in the richness parameter.

\subsection{Depth and the comparison of architectures}

Everything above enters an architecture through $F_0(\ell)$, one Krylov series
on the initial kernel and error per design, and through $\bar H_2$. For two
designs on the same dataset, the comparison principle of Corollary~6 applies
with the bounds of this section.

\begin{corollary}[Comparing designs on the rough part]
\label{cor:depth}
Let designs $A$ and $B$ be trained by gradient flow on the same dataset and
let $T>0$.
\begin{enumerate}
\item[(i)] (A priori.) $\|E^A(T)\|$ is at least the right side of
Eq.~\eqref{eq:rough-floor-apriori} computed with $F_0^A$ and $\bar H_2^A$ at
$\ell=\sqrt{T/2}\,\|E_0^A\|$, and at least $\tfrac12\rho^A(\gamma^\star_A(T))$.
\item[(ii)] (Along the run of $B$, Section~\ref{sec:chain}.) $\|E^B(T)\|^2$ is at most the
ceiling of the cutoff bound computed for $B$: the energy of its initial error on
the modes that are still unfit at $T$ under its lifted spectrum, plus the
cross-front leakage and the decayed fit part. Its inputs (the transported
unfit set, the leakage, the fit residual) are read along $B$'s trajectory;
it is not computable before training.
\item[(iii)] If the bound in (ii) is below the bound in (i), then
$\|E^B(T)\|<\|E^A(T)\|$: $B$ has learned more of the rough part of the
labels by time $T$. Since (ii) needs $B$'s run, the corollary compares a
trained $B$ against an untrained $A$; a before-training comparison of two
designs needs (ii) replaced by an a~priori ceiling, which the theory does not
have.
\end{enumerate}
\end{corollary}

The corollary separates what depth does into two effects that can be checked
one at a time. A deeper network can have a larger initial kernel on the slow
part of the labels, a smaller $\rho(\gamma)$ at every $\gamma$; that lowers
the held rough error after the lazy phase, since the cutoff of
Definition~\ref{def:cutoff} starts lower. And it can have a larger
$\bar H_2$, the rate at which it can move its kernel per unit of path; that
shortens the ignition time of Corollary~\ref{cor:min-time-rough} and advances
the onset of the lift given by the richness expansion of
Section~\ref{sec:chain}. Both are computed before training, one Krylov series
and one nested automatic differentiation per design; the theory says which
part of the error each one removes and in what time, but not that either
grows with depth, which is a measurement or a statement about the
initialization.

\subsection{Scope and what to measure}

\textbf{Exact under gradient flow, no hypothesis beyond smoothness.}
Proposition~\ref{prop:rough-energy}; Theorem~\ref{thm:rough-apriori} and
Corollary~\ref{cor:min-time-rough} with $\bar H_2$ taken over the ball of
radius $\ell(t)$; Proposition~\ref{prop:cutoff-floor}; Corollary~\ref{cor:depth}(i).
\textbf{Proved in Section~\ref{sec:chain}, with named inputs read along the run.} The cutoff bound (two-sided, with
the leakage as remainder, the leakage not yet evaluated at any sample size);
Corollary~\ref{cor:depth}(ii)--(iii) as a trajectory bound; the onset law
inside its radius.
Nothing in this section is approximated; the closure hypothesis of the earlier
draft is not needed, and its failure was the failure of a fixed projector.

\textbf{What the a priori floor is worth along a run [n=200, 36 runs, to be
redone].} Evaluated on the saved kernels and residuals of the ground-truth runs
with the Krylov functionals of $(K_0,E_0)$:
Eq.~\eqref{eq:rough-floor-apriori} is halved by its $\bar H_2$ term at a path
length of about $0.6$, which gradient flow reaches at $t\approx1$ of $3000$,
and is negative from $t\approx10$, because $\bar H_2$ is the worst-case
curvature and the realised drift of the kernel is far smaller; the first-order
term $F_0(\ell)$ alone exceeds the measured error at every later time, by
$2\times$ at $t=100$, $4\times$ at $t=300$ and $17\times$ at $T$ (medians), so
the whole content of the late error is in the motion of the kernel. The path
floor with the actual kernel along the run, charged by the measured cost of a
functional fixed at the start, holds $0.94$ of the error at $t=10$ and is
empty from $t\approx100$, the cost along that functional having been lifted
to $30\times$ its first-order value by then and $280\times$ by $T$: a floor
charged by the lift loses content exactly when the lift arrives. The energy
form of the earlier draft is empty from the start for these functionals, its
coupling term being $10^3$ times its main term at $t=10$. The one-time floor
at the current state, $\|E_{1/t}(t)\|$, holds $0.92$ of the error at $t=10$ and
$0.70$--$0.74$ from $t=100$ on, and is the floor that stays informative.
(Numbers: EXPERIMENTS.md, item C1.)
The measurements that decide whether the closure is the right one: (13) the
cost $\int\|\varphi\|_{K(s)}|\dot\theta|$ of the Krylov functionals along each
run against its a~priori bound $\ell\|\varphi\|_{K_0}+\tfrac12\bar H_2\ell^2$,
which is the question of how loose $\bar H_2$ is; (14) the slow part
$\|E_\gamma(t)\|$ at the current state against $\|E(t)\|$, the one-time floor
that is exact to first order; (15) $\rho(\gamma)$ for the labels against the
measured error's spectrum in the current kernel basis; (16) the cutoff rate
$\gamma^\star(t)$ against the rate at which the measured error's spectrum
drops, which tests Proposition~\ref{prop:cutoff-floor} as an estimate rather
than as a bound; (17) $\rho(\gamma)$ and $\bar H_2$ as functions of depth at
initialization, which are the quantities through which depth enters
Corollary~\ref{cor:depth}.
